\begin{abstract}
Statistical watermarks such as KGW bias a language model toward a keyed ``green'' subset of the vocabulary so that its outputs can later be attributed. Deployments treat the secret key as a free parameter: any valid key is assumed to be behaviorally interchangeable. For code generation this assumption is consequential, because a single token can separate a secure program from a vulnerable or broken one, yet evaluations of watermarked code use one key and measure correctness, not security. We formulate the question as two distinct estimands, the average effect of watermarking relative to matched unwatermarked generation and the between-key heterogeneity among valid keys, and give a paired, clustered measurement protocol with fail-closed manipulation gates. We show that between-key contrasts do not depend on the shared baseline, and derive a first-order output-length requirement for per-output detection. We then report pilot evidence for one open code model under a preregistered setting. A calibration of 7,200 generations shows a monotone dose response in detector strength. A ten-key pilot of 792 programs shows functional pass rates ranging from 12.5\% to 70.8\% across keys, while security labels are nearly key-invariant and only 2.9\% of watermarked outputs are individually detectable. We therefore do not claim a key-dependent security effect; we identify the validation gates that must pass before such a claim can be tested.
\end{abstract}
